\chapter{A Solver Architecture for Reviewable Candidates}

\section{Architecture begins at a responsibility boundary}

\paragraph{Reader bridge.}
The previous chapter produced a candidate contract. Architecture is needed
because that contract can be lost when parsing, solving, visualization, and
model editing are treated as one operation.

A solver service may consume a read-only problem declaration and return
candidate evidence. It must not infer the engineering purpose, broaden the
released freedoms, replace source observations, decide applicability, or
write a result into an authoritative Alignment merely because a numerical run
terminated.

For the running discrepancy, the conceptual flow is
\[
\begin{array}{c}
\boxed{\mathcal P}\longrightarrow\boxed{\text{validate and assemble}}\\[0.4em]
\downarrow\\[-0.2em]
\boxed{\text{solve and verify}}\longrightarrow\boxed{\mathcal C}\\[0.4em]
\downarrow\\[-0.2em]
\boxed{\text{evaluate and decide}}
\end{array}
\]
Only the middle operations belong to the solver responsibility. The final
arrow is an engineering workflow boundary.

\section{The problem package}

A portable problem package needs more than geometry:
\[
\mathcal P=
(id,I,P,M,x_0,\mathcal X_{\mathrm{free}},q,\mathcal F,
J,W,\Sigma,\pi),
\]
where \(W\) records scaling and \(\pi\) records provenance. Validation checks
that every free variable is declared, fixed inputs remain locked, units are
compatible, objective terms name their source, and all referenced models and
contexts are available. Failure of this check is not solver non-convergence;
it is an incomplete or inapplicable problem declaration.

\section{Assembly makes mathematical structure explicit}

\paragraph{Reader bridge.}
Once the engineering contract is complete, the architecture can translate it
into numerical objects without changing its meaning.

Assembly constructs:
\begin{itemize}
\item the parameter vector and bounds from the released freedoms;
\item equality and inequality functions from declared constraints;
\item separately typed objective and diagnostic residuals;
\item reconstruction and observation evaluators;
\item scaling, derivative providers, and sparsity patterns;
\item termination tolerances and independent verification settings.
\end{itemize}

Residuals are not interchangeable merely because they are stored in one
vector. Type and provenance must survive assembly so that a reviewer can
distinguish an endpoint constraint, an observation mismatch, a regularizer,
and a numerical diagnostic.

\section{Why SQP, derivatives, and sparsity belong behind interfaces}

\paragraph{Reader bridge.}
SQP is useful for a smooth nonlinear problem with visible equality and
inequality constraints. It is one replaceable numerical method, not the
owner of the problem.

At iterate \(x_k\), an SQP implementation forms a local quadratic model and
linearized constraints, obtains a step \(p_k\), and applies a globalization
rule such as a line search or trust region. The architecture must expose:
\[
x_k,\quad J(x_k),\quad r(x_k),\quad
\nabla J(x_k),\quad A(x_k),\quad p_k,
\]
together with constraint margins, derivative provenance, scaling, and solver
diagnostics.

Analytic, automatic, or finite-difference derivatives may be supplied if their
method and validation are recorded. Sparsity is a computational property of
the declared dependency graph. Exploiting it can reduce work; it does not
justify an undeclared objective or constraint.

\section{Termination, convergence, and verification}

A solver may stop because a tolerance was met, progress stalled, an iteration
limit was reached, or evaluation failed. These outcomes must remain distinct.
A nominal convergence flag is credible only with the corresponding evidence,
including:
\begin{itemize}
\item objective and step histories;
\item equality residuals and inequality margins;
\item a stationarity or optimality measure appropriate to the method;
\item derivative and scaling checks where applicable;
\item reconstruction on the declared verification grid;
\item warnings, active assumptions, and numerical error estimates.
\end{itemize}

This evidence establishes what happened within \(M\) and \(\mathcal P\). It
does not establish that the observation model is correct, that the candidate
is constructible, or that the purpose warrants adoption.

\section{The candidate package and review boundary}

The service returns, without applying,
\[
\mathcal C=
(\mathcal P.id,M.id,x^\star,\widehat\gamma^\star,
r^\star,J^\star,\text{margins},\text{termination},
\text{verification},\Sigma,\pi).
\]
A review surface may compare candidates, inspect residuals and sensitivities,
and preview a representation. Preview is not Physical Realization; comparison
is not evaluation under every relevant purpose; an ``Apply'' operation, if a
future workflow provides one, requires a separate authorized command and
preserved provenance.

\section{Current AXTRAN evidence and the general architecture}

The architecture above is a scientific responsibility model and a direction
for implementation. Current repository evidence for AXTRAN is narrower:
AXTRAN is experimental; a single SQP step and a finite repeated loop are
evidence of a numerical path. They do not demonstrate a general optimizer
architecture, production robustness, live optimization, network
optimization, vehicle-aware objectives, cant optimization, or automatic model
integration.

Accordingly, examples of richer sessions, visualization, alternative solvers,
or network assembly are future or problem-specific possibilities. They must
not be read as present AIM or AXTRAN capability.

\section{What the engineer receives}

For the CAD/survey discrepancy, the architecture returns one or more
traceable curvature candidates and the evidence needed to question them:
which parameters moved, what stayed fixed, which residuals improved, which
constraints are close to active, why the run stopped, and which uncertainties
remain. Applicability is then checked under the intended project and
realization context. Evaluation compares consequences under declared
purposes. An authorized Engineering Decision may accept, reject, or request a
new problem formulation.

\begin{summary}
A solver architecture is trustworthy when it preserves the engineering
problem on the way in and returns a reviewable candidate on the way out.

\[
\boxed{\text{calculated}\neq\text{applicable}\neq
\text{accepted}\neq\text{authoritative}}
\]

The solver owns numerical work inside the declared contract. It does not own
Alignment identity, applicability, acceptance, or authority.
\end{summary}
